\documentclass{article}
\usepackage{tikz}
\usetikzlibrary{angles, quotes}
\usepackage{amssymb}
\usepackage{pgfplots}
\pgfplotsset{width=10cm,compat=1.9}
\usepgfplotslibrary{statistics}
\usepackage[left=2.54cm,right=2.54cm,top=2.54cm,bottom=2.54cm]{geometry}
\usepackage{tasks} %Paquete necesario para la producción de items horizontales para algunos ejercicios de tarea empleando el comando /begin{multicols}}{#}.../end{multicols} para ayudarnos a reducir las páginas
\usepackage{pdflscape} %comando que permite cambiar a orientación horizontal las hojas%
\usepackage{soul} %Paquete para permitir la compilación de acentos usando el comando {\hl{}}}\\
\sethlcolor{yellow(munsell)} %comando para definir el color para el subrayado usando el comando \hl
\usepackage[utf8]{inputenc}
\usepackage[spanish]{babel}
\usepackage{icomma}
\usepackage{ragged2e}
\usepackage{siunitx}
\usepackage{url}
\usepackage[colorlinks=true, urlcolor=blue,  linkcolor=black, citecolor=green]{hyperref}
\usepackage{pdfpages}
\usepackage{blindtext} %Paquete que genera texto ficticio (dummy text) con el comando: \blindtext
\usepackage{cancel} %in the preamble gives you four different modes of striking through: \cancel{text to cancel} draws a diagonal line (slash) through its argument, \bcancel{text to cancel} uses the negative slope (a backslash), \xcancel{text to cancel} draws an X (actually \cancel plus \bcancel), \cancelto{〈value〉}{〈expression〉} draws a diagonal arrow through the 〈expression〉pointing to the 〈value〉 (math-mode only)
\usepackage{csquotes}
\usepackage{afterpage}
\usepackage{parskip} 
\usepackage{float}
\usepackage{enumitem}
\usepackage{multicol}%Paquete que permite la creación aislada de columnas de texto. De esta forma se reduce la cantidad de páginas en nuestro documento
\usepackage{lipsum} %Paquete que genera texto de relleno al igual que \usepackage{blindtext}, se usa con el comando \lipsum[#-#] los #(hashtags) delimitan el número de páginas que deseamos ocupar del paquete lipsum para usarlos como relleno. 

% Definir el comando personalizado para la nota
\newcommand{\mynote}[1]{%
\begin{tikzpicture}[baseline=-0.75ex]
\node[draw, circle, fill=Apple Green!60, inner sep=2pt] (note) {\textbf{Nota:}};
\end{tikzpicture}%
\ \textit{#1}%
}

\newenvironment{Figura}
{\par\medskip\noindent\minipage{\linewidth}}
{\endminipage\par\medskip}
\usepackage{caption}
\usepackage[
backend=biber,
sorting=none,
url=true
]{biblatex} %bibliografía
\addbibresource{biblio.bib}
% Establecer el espacio entre las entradas de la bibliografía
\setlength{\bibitemsep}{1\baselineskip} % Puedes ajustar el valor según tus preferencias
\usepackage{amsmath,amsthm,amssymb,amsfonts}
\usepackage{pifont} %Permite la compilación del comando \xmark: es el símbolo de la tachita
\usepackage{mathtools} %permite la compilación de símbolos para matrices
\usepackage{empheq} %Paquete que se relaciona con \usepackage[most]{tcolorbox} para la creación de cuadros/cajas de colores para delimitar los resultados matemáticos o líneas de texto usando la línea de comando: \begin{empheq}[box={\mymath[colback=orange(webcolor)!70,drop lifted shadow]}]{equation*} \end{empheq}
\usepackage[most]{tcolorbox} %Paquete que se relaciona con \usepackage{empheq} para la creación de cuadros/cajas de colores para delimitar los resultados matemáticos o líneas de texto
\renewcommand{\qedsymbol}{$\blacksquare$}
\usetikzlibrary{positioning,decorations.pathreplacing} %Paquete que permite la compilación de llaves para cuadros sinópticos (brace diagram) 
\usepackage{schemata} %The schemata package is designed just for these brace diagrams
\usepackage{booktabs}

\newcommand\AB[2]{\schema{\schemabox{#1}}{\schemabox{#2}}} %comando o paquete necesario para crear cuadros sinópticos

\newtcbox{\mymath}[1][]{%
nobeforeafter, math upper, tcbox raise base,
enhanced, colframe=blue!30!black,
colback=blue!30, boxrule=1pt,
#1} %Comando importante para encasillar los resultados matemáticos en cajas de diferentes colores, se relaciona con los paquetes \usepackage{empheq} y \usepackage[most]{tcolorbox} 

\newenvironment{sysmatrix}[1]
{\left(\begin{array}{@{}#1@{}}}
{\end{array}\right)}
\newcommand{\ro}[1]{%
\xrightarrow{\mathmakebox[\rowidth]{#1}}%
}
\newlength{\rowidth}% row operation width
\AtBeginDocument{\setlength{\rowidth}{3em}}  %Comando importante para laproducción de lineas de operaciones entre matrices, método de Gauss o Gauss Jordan se relaciona con \newenvironment{sysmatrix}

%formato para cambiar el horario a español
\usepackage[spanish]{datetime2}
\DTMsetdatestyle{spanish}

\renewcommand{\today}{\DTMdisplaydate{\the\year}{\the\month}{\the\day }{-1}}
%%%%%%%%%%%%%%%%%%%%%%%%%%%%%%%%%%%%%%%%%

\usepackage{datetime}
\newcommand{\mycurrenttime}{\xxivtime}

%%%%%%%%%%%%%%%%%%%%%%%%%%%%%%%%%%%%%%%%%%%%%%%%%%%%%%%%%%%%%%%%%%%%%%%%%%%%%%%%%%%%%%%%%%%%%%%%%%%%%%%%%%%%%%%%%%%%%%%%%%%%%%%%%%%%%%%%%%%%%%%%%


%%%%%%%%%%%%%Caja de color para el título%%%%%%%%%%%%%%%%%%%%%%%%%%%%%%%%%%%%%%%%%%%%%%%%%%%%%%

\definecolor{myframecolor}{RGB}{1, 45, 75} % Prussian Blue
\definecolor{myboxcolor}{RGB}{0, 107, 92} % Apple Green

\newtcolorbox{mybox}{
enhanced,
colback=myboxcolor!7, % Color de fondo del cuadro
colframe=myframecolor, % Color del marco del cuadro
arc=0pt,
boxrule=1pt,
borderline west={2mm}{-10mm}{myframecolor}, % Borde en el lado izquierdo
sharp corners=southwest,
width=\linewidth,
before=\par\vspace{\bigskipamount}, % Espacio antes del cuadro
after=\par\vspace{\bigskipamount} % Espacio después del cuadro
}


\newcommand{\euler}{e} %Comando para producir letra e de euler
\newenvironment{remark}{\par\vfill\footnotesize % Vertical white space above the remark and smaller font size
\begin{list}{}{
	\leftmargin=80pt % Indentation on the left
	\rightmargin=60pt}\item\ignorespaces % Indentation on the right
\makebox[-2.5pt]{\begin{tikzpicture}[overlay]
		\node[draw=Horizon!60,line width=2.5pt,circle,fill=Horizon!25,font=\sffamily\bfseries,inner sep=4pt,outer sep=0pt] at (-19pt,5pt){\textcolor{Horizon}{Nota}};\end{tikzpicture}} % Blue Nota in a circle
\advance\baselineskip -1pt}{\end{list}\vskip5pt} % Tighter line spacing and white space after remark
\usepackage{graphicx}

\usepackage{titling}

\input{colores.tex}

%%Producir signos de cita textual``Asignatura''

\newcommand{\xmark}{\ding{55}} %%%comando que genera la tachita

\newenvironment{MyColorPar}[1]{%
\leavevmode\color{#1}\ignorespaces%
}{%
}%

%%%%%%%%%%%%%%%%%%% Título de la tarea, Nombre de alumno

\title{Clase \# 39 - Mecánica Estadística}
\author{\emph{Notas de clase} $\boldsymbol{\mid}$}
\date{\today}

%%%%%%%%%%%%%%%%%%% Datos de la Materia

\usepackage{fancyhdr}
\fancypagestyle{plain}{%  the preset of fancyhdr 
\fancyhf{} % clear all header and footer fields
\fancyfoot[R]{
\includegraphics[width=2cm]{LogoFCFMBUAP (1).png}}
\fancyfoot[L]{{\bfseries{\thedate{}}} a las {\bfseries{\mycurrenttime{}}} horas{} (GMT-6, H. Puebla de Zaragoza, Pue)}
\fancyhead[L]{Mecánica Estadística}
\fancyhead[R]{\theauthor}
}
\makeatletter
\def\@maketitle{%
\newpage
\null
\vskip 1em%
\begin{center}%
\let \footnote \thanks
{\LARGE \@title \par}%
\vskip 1em%
%{\large \@date}%
\end{center}%
\par
\vskip 1em}
\makeatother

\usepackage{lipsum}  
\usepackage{cmbright}


%%%%%%%%%%%%%%%%%%%%%%%%%%%%%%%%%%%%%%%%%%%%%%%%%%%%%%%%%%%%%%%%%%%%%%%%%%%%%%%%%%%%%%%%%%%%%%%%%%%%%%%%%%%%%%%%%%%%%%%%%%%%%%%%%%%%%%%%%%%%%%%%%%%%%%%%%%%%%%%%%%%%%%%%%%%%%%%%%%%%%%%%%%%%%%%%%%%%%%

\begin{document}



\maketitle


%%%%%%%%%%%%%%%%%%% Datos del profesor, campus y aula

\begin{mybox}
\noindent\begin{tabular}{@{}ll}
	{\bfseries{Profesor}} & Dr. Jorge Velázquez Castro\\
	\textcolor{prussianblue}{\bfseries{Campus}} / \textcolor{trueblue}{\bfseries{Facultad}}  & \textcolor{prussianblue}{\bfseries{C.U. BUAP}} /  \textcolor{trueblue}{\bfseries{FCFM}} \\
	\textcolor{prussianblue}{\bfseries{Edificio}} /  \textcolor{trueblue}{\bfseries{Salón}}    & 	\textcolor{prussianblue}{\bfseries{EMA4}} /  \textcolor{trueblue}{\bfseries{407}}   
\end{tabular} 
\end{mybox} \vspace{1cm}

\tableofcontents

\vspace{2cm}

\section{Clase de hoy: Ensamble gran canónico}

\begin{figure}[H]
\centering
    \includegraphics[width=0.8\textwidth]{2.png}
\caption*{.}
\end{figure}


\begin{figure}[H]
\centering
    \includegraphics[width=1.1\textwidth]{1.png}
\caption*{.}
\end{figure}


En términos del espacio fase clásico (continuo).

Tenemos:

\[ \rho_{gc} (N_{s}, q^{\prime})_{v}, p^{\prime}_{v} ) dq^{\prime}_{v} dp^{\prime}_{v} =  \frac{ e^{ - \beta (H (p^{\prime}_{v}, q^{\prime}_{v}) + \mu_{s} N_{s} ) } }{\Theta (\beta, v,\mu)  } dp^{\prime}_{v} dq^{\prime}_{v}  \]


Donde la función gran canónica se calcula así:

\[ \Theta_{Distinguibles} = \sum\limits_{n=0}^{\infty} \frac{1}{h^{3N}} \int\limits dp^{3N} dq^{3N} e^{- \beta (  H_{N} (q^{\prime}_{s}, p^{\prime}_{s}  )  - N \mu  )  } \]

Tenemos:


\[  \Theta_{Indistinguibles} = \sum\limits_{N=0}^{\infty} \frac{1}{N!h^{3N}} \int\limits dp^{3N} dq^{3N} e^{-\beta  (H_{N} (q^{\prime}_{s},  p^{\prime}_{s} ) - \mu N )  }    \]



Para relacionarlo con cantidades termodinámicas calculamos:



\[S = \langle -K_{B}\ln{ \rho_{gc} } \rangle  \]

Tenemos:


\[ S = \langle -K_{B} \left\{   \ln e^{ -\beta (H - \mu N_{s})  } - \ln \Theta \right\} \rangle_{gc} \quad\text{es decir, es promedio respecto al gran canónico (gc)}  \]


Tenemos

\[  \langle \frac{1}{T} H \rangle - \langle \frac{\mu}{T} N_{s} \rangle + K_{B}\ln \Theta   \]


Tenemos:


\[ S = \frac{1}{T} \langle \underbrace{H}_{U} \rangle - \frac{\mu}{T} \langle N_{s} \rangle + K_{B} \ln \Theta    \]


\[ \Rightarrow K_{B}\ln \Theta (\frac{1}{T}, V, \frac{\mu}{T}  ) = S - \frac{1}{T}U + \frac{\mu}{T} \langle N_{s}\rangle  \equiv S \left[ \frac{1}{T}, \frac{\mu}{T}  \right] \]



O bien en representación energética:



\[  -K_{B}T \ln \Theta (T,V,\mu) = U-TS - \mu \langle N_{S} \rangle \equiv \underbrace{\Phi (T,V,\mu)}_{\text{Potencial gran canónico}}  \]

Notar que:



\[ \Theta(T,V,\mu) = \sum\limits_{N=0}^{\infty} e^{ \frac{\mu  }{K_{B}T} N } Z(T,V,N)  \]

Todavía más; para partículas independientes:



\[ Z(T,V,N ) = \frac{1}{N!}  \left[ Z(T,V,1)   \right]^{N}  \]


\[ \Theta (T,V,\mu) = \sum\limits_{N=0}^{\infty} \frac{1}{N!} \left[  e^{\frac{\mu}{K_{B}T}  } Z(T,V,1)     \right]^{N}   \]


\[= e^{ \left\{  e^{  \frac{ \mu   }{K_{B}T}       } Z(T,V,1)        \right\}     }  \]

Nota:


AL término \[e^{\mu/k_{B}T}\] se le llama fugacidad
 
 \subsection{Ejemplo de otro potencial y su función de partición }



\[  \gamma\left(  \frac{1}{T}, \frac{P}{T}, N  \right) = S - \frac{1}{T}U - \frac{P}{T}V \quad \Rightarrow \quad -T\gamma = U - TS +PV \]


Tenemos:


\[  \bar{\gamma} = \int\limits \sum\limits_{i} e^{  - \frac{1  }{ K_{B}T } U_{i} - \frac{P}{K_{B}T} V  } \Omega(U_{i},V,N) dV    \]



Tenemos:

Si $\left. K_{B}\ln{\bar{\gamma}} = \gamma \right\} \text{Potencial en la entrópica}  $


Entonces:


\[ \underbrace{-K_{B}T\ln{\bar{\gamma}} = \varphi}_{\text{Energética}}  \]

\subsection{Gas ideal (Gran canónico)}

Sabemos que:


\[ Z(T,V,1) = \frac{V}{\lambda^{3}} \quad \text{con}  \lambda = \left( \frac{ h^{2} }{2\pi m K_{B}T}    \right)^{\frac{1}{2}}  \]



También sabemso que para particulas independientes e indistinguibles:


\[    \Theta = \sum\limits_{N=0}^{\infty} \frac{1}{N!} e^{ \frac{ \mu  }{ K_{B}T }N   } \left[  Z(T,V,1)   \right]^{N}       \]

Por lo tanto:

\[ \Theta = exp \left\{ e^{\frac{ \mu }{K_{B}T  }} \frac{ V  }{\lambda^{3}}    \right\} = exp\left\{  e^{  \frac{ \mu }{K_{B}T  }  } V \left(  \frac{2\pi K_{B}T  }{h^{2}}  \right)^{\frac{3}{2}} \right\}    \]



El potencial gran canónico es:




\[ Phi(T,V,\mu) = -K_{B}T \ln \Theta = -K_{B}Te^{ \frac{ \mu }{K_{B}T}  } V \left( \frac{2\pi m K_{B}T}{h^{2}}   \right)^{ \frac{3}{2} } \]





Las ecuaciones de estado son:



\[   - \frac{  \partial \phi }{ \partial T }  = S(T,V,\mu) = e^{\frac{\mu}{K_{B}T}  } VK_{B} \left(  \frac{ 2\pi mK_{B}T }{h^{2}}    \right)^{ \frac{3}{2} } \left(   \frac{5}{2} - \frac{\mu}{K_{B}T}    \right)   \]



Tenemos:


\[   - \frac{  \partial \phi  }{\partial V} = P(T,V, \mu) = K_{B}T e^{ \frac{ \mu  }{K_{B}T}  }  \left( \frac{ 2\pi mK_{B}T }{h^{2} }  \right)^{\frac{3}{2}}      \]


\textbf{NOTA}:


Como:



\[  -K_{B}T\ln\Theta = U -TS - \mu{N}  \]

Usando la relación de Euler ($U=TS-PV+\mu{N} $)


\[ -K_{B}T \ln \Theta = ( TS-PV+\mu N ) -TS -\mu N   \]



\[ \Rightarrow  \quad \boxed{  \ln{\Theta} )= \frac{PV}{K_{B}T}    }  \]


Por eso alguns veces se le llama el \underline{potencial PV}


\subsection{ Fluctuaciones del gran canónico  }


Tenemos:



\[ \rho_{i,N}  =  \frac{ e^{ -\beta (E_{i} - \mu{N}) }  }{  \Theta }   \]


 O bien:
 
 
 
 \[  \rho_{E,N} = \frac{ \Omega_{N}(E) e^{  -\beta (  e-\mu{N} ) }   }{  \Theta }         \]
 
 
 
 El número más probable de partículas, es el está en el máxima de $P_{E,N}$ 



Para encontrarlo derivamos e igualamos a cero.



\[  \frac{\partial }{\partial N}  \ln \rho_{E,N} = \frac{  \partial  }{\partial N}  \left[  \ln \Omega_{N}(\epsilon) + \frac{\mu}{K_{B}T}N  \right]_{ N=N^{\ast} } =0 \]





\[  \Rightarrow \quad \frac{  \partial }{\partial N} K_{B}\ln \Omega_{N^{\ast}}(E) = - \frac{\mu}{T}   \]



\[  \Rightarrow \quad  \underbrace{ \frac{\partial }{\partial N}   S(E,V,N^{\ast}) = - \frac{\mu}{T}}_{\text{Ec. de estado que me indica el valor más probable de $N$ en $N^{\ast}$ }}  \]

Por otro lado, podemos calcular el promedio:



\[  \langle N \rangle = \sum\limits_{i,N} N \rho_{i,N} = \sum\limits_{i,N}  \frac{ N }{\Theta} e^{  -\beta (E_{i}  - \mu N )  } = \frac{\partial }{\partial \mu} K_{B}T \ln \Theta = - \frac{  \partial \Phi }{\partial \mu } \]


Es decir:



\[   \underbrace{-\frac{ \partial \Phi   }{\partial \mu} = \langle N \rangle}_{ \text{Ec. de estado en Gran Canónico} }   \]

 
\end{document}


